\clearpage
\fancyhead[L]{\small VERSCHIEDENE AUFSÄTZE ÜBER ASTRONOMIE}
\SourceLeaf{78}%
\pdfbookmark[1]{Berechnung des Osterfestes}{osterfest1800}
\begin{center}
{\large BERECHNUNG DES OSTERFESTES.\par}
\medskip\hrule\medskip
{\small Monatliche Correspondenz zur Beförderung der Erd- und Himmels-Kunde,\par
herausgegeben vom Freiherrn \textsc{von Zach}. August 1800.\par}
\medskip\hrule
\end{center}
\medskip
Die Absicht dieses Aufsatzes ist nicht, das gewöhnliche Verfahren zur Bestimmung
des Osterfestes zu erörtern, das man in jeder Anweisung zur mathematischen
Chronologie findet, und das auch an sich leicht genug ist, wenn man
einmal die Bedeutung und den Gebrauch der dabei üblichen Kunstwörter, \textit{güldne
Zahl}, \textit{Epacte}, \textit{Ostergränze}, \textit{Sonnenzirkel} und \textit{Sonntagsbuchstaben} weiss, und die
nöthigen Hülfstafeln vor sich hat: sondern von dieser Aufgabe eine von jenen
Hülfsbegriffen unabhängige und blos auf den einfachsten Rechnungs-Operationen
beruhende rein analytische Auflösung zu geben. Hoffentlich wird dieselbe
nicht allein dem blossen Liebhaber, dem jene Methode nicht geläufig ist, oder
der wohl in den Fall kommt, die Bestimmung der Zeit des Osterfestes unter Umständen,
wo ihm die nöthigen Hülfsmittel nicht zur Hand sind, oder für ein Jahr,
worüber er keinen Kalender nachschlagen kann, auf der Stelle zu wünschen,
nicht unangenehm sein, sondern sich auch dem Kenner durch ihre Einfachheit
und Geschmeidigkeit empfehlen. Die folgenden Vorschriften, die jeder, der es
der Mühe werth hält, leicht wird ins Gedächtniss fassen können, gelten für zwei
Jahrhunderte, von 1700 bis 1899; sie können aber auch leicht, durch gehörige
Veränderung der darin vorkommenden beständigen Zahlen und mit Beifügung
einer unerheblichen Ausnahme, die eine Folge der Einrichtung unsers Kalenders
ist, und zufälliger Weise während jenes Zeitraumes nicht Statt findet, für jedes
andere gegebene Jahrhundert eingerichtet werden.

\SourceLeaf{79}%
\expandafter\def\csname HistoricalNoteText2\endcsname{Der Grund davon liegt darin, dass $a$ nur 19 verschiedene Werthe $(0,\ 1,\ 2\ldots18)$ bekommen
kann, und folglich auch $d$ nur eben so viele, unter welchen der Werth 29 nicht mit begriffen ist.}
I. Man dividire die Zahl des Jahres, für welches man Ostern berechnen
will, mit 19, mit 4 und mit 7, und nenne die Reste aus diesen Divisionen, respective
$a$, $b$ und $c$. Geht eine Division auf, so setzt man den zugehörigen Rest
$=0$; auf die Quotienten wird gar keine Rücksicht genommen. Eben das gilt
von den folgenden Divisionen.

II. Man dividire ferner $19a+23$ mit 30, und nenne den Rest $d$.

III. Endlich dividire man $2b+4c+6d+3$, oder $2b+4c+6d+4$, je
nachdem das vorgegebene Jahr zwischen 1700 und 1799, oder zwischen 1800
und 1899 inclus. liegt, mit 7, und nenne den Rest $e$.

Alsdann fällt Ostern auf den $22+d+e^{\mathrm{ten}}$ März, oder wenn $d+e$ grösser
als 9 ist, auf den $d+e-9^{\mathrm{ten}}$ April.

\textit{Beispiele.} Für das Jahr 1744 findet man bei der Division der Zahl 1744
mit 19 den Rest $15=a$; die Division mit 4 geht auf, also $b=0$; die Division
mit 7 gibt den Rest $1=c$. Hieraus wird $19a+23=308$, welches mit
30 dividirt den Rest $8=d$ gibt. Endlich gibt $2b+4c+6d+3=55$ mit 7
dividirt den Rest $6=e$. Folglich ist Ostern den $22+8+6^{\mathrm{ten}}$ März, oder den
$14-9$ d.\,i. den 5. April.

Für 1800 wird $a=14$, $b=0$, $c=1$; $19a+23=289$, also $d=19$;
$2b+4c+6d+4=122$, also $e=3$; mithin Ostern den $19+3-9$ d.\,i. den
13. April.

Für 1818 ist $a=13$, $b=2$, $c=5$; $19a+23=270$, also $d=0$;
$2b+4c+6d+4=28$, also $e=0$, folglich Ostern den 22. März.

In dem letzten Beispiele fällt Ostern auf den möglich frühesten Tag, denn
es ist einleuchtend. dass $d$ und $e$ hier ihre möglich kleinsten Werthe haben.
Von der andern Seite erhellt, dass Ostern nie später als den $22+29+6^{\mathrm{ten}}$ März,
d.\,i. den 26. April eintreten könne, da $d$ nicht grösser als 29, und $e$, nicht
grösser als 6 werden kann; allein in dem achtzehnten und neunzehnten Jahrhundert
kann nie $d=29$ werden\HistoricalNote{2}; der späteste Ostertag ist folglich, während
dieses Zeitraumes, der 25. April, welcher Statt hat, wenn zugleich $d=28$ und
$e=6$ wird. Diese beiden Bedingungen vereinigen sich in den Jahren 1734
und 1886. In andern Jahrhunderten könnte zwar $d=29$ werden, allein ge%
\SourceLeaf{80}%
rade in diesem Falle tritt die oben erwähnte Ausnahme ein, vermöge welcher alsdann
der Werth von $d$ wieder auf 28 heruntergebracht wird, so dass der 25.
April der absolut späteste Ostertag ist. Eine weitere Entwickelung dieses Umstandes
würde hier zu weitläuftig werden.

Die Analyse, vermittelst welcher obige Formel gefunden wird, beruhet eigentlich
auf Gründen der \textit{höhern Arithmetik}, in Rücksicht auf welche ich mich
gegenwärtig noch auf keine Schrift beziehen kann, und lässt sich daher freilich
in ihrer ganzen Einfachheit hier nicht darstellen: inzwischen wird doch folgendes
hinreichen, um sich von dem Grunde der Vorschriften einen Begriff zu machen
und von ihrer Richtigkeit zu überzeugen.

I. Die güldene Zahl eines Jahres unserer Zeitrechnung ist bekanntlich der
Rest, der entsteht. wenn man zu der Jahres-Zahl 1 addirt und die Summe mit 19
dividirt; nur muss derselbe $=19$ gesetzt werden, wenn die Division aufgeht.
Daraus folgt leicht, dass $a+1$ die güldene Zahl des vorgegebnen Jahres sein werde.

II. Die Oster-Gränze, das ist der Tag des Oster-Vollmonds, fällt im 18.
und 19. Jahrhundert für ein Jahr, dessen güldne Zahl 1 ist, auf den 13. April,
und alsdann den ganzen Zirkel von 19 Jahren hindurch, d.\,i. bis zum Jahre,
dessen güldne Zahl 19 ist, inclus., in jedem Jahre \textit{entweder} 11 Tage früher, \textit{oder}
19 Tage später, als in dem nächst vorhergehenden, je nachdem sie in diesem
\textit{entweder} in den April \textit{oder} in den März gefallen war, wie man sich leicht aus einer
Tafel der Oster-Gränzen überzeugen kann; folglich in dem Jahre, dessen
güldne Zahl 2 ist, auf den 2. April, in dem folgenden auf den 22. März, in dem
Jahre, dessen güldne Zahl 4 ist, auf den 10. April u.\,s.\,f. Hieraus folgt, dass
die Oster-Gränze nie \textit{vor} den 21. März und nie \textit{nach} dem 19. April fällt; nimmt
man also an, sie falle für das Jahr, dessen güldne Zahl $a+1$ ist, auf den
$21+D^{\mathrm{ten}}$ März (indem man die Tage des Aprils auf den März reducirt), so liegt
$D$ allemal zwischen Gränzen 0 und 29 inclus. Für $a=0$ ist also $D=23$,
für $a=1$ wird $D=23-11$, für $a=2$ wird $D=23-2\times11$, für $a=3$
wird $D=23-2\times11+19$ u.\,s.\,f.; und allgemein $D=23-11p+19q$, wo $p$
und $q$ durch die Bedingungen bestimmt werden, dass $p+q=a$ werde und $D$
zwischen die Gränzen 0 und 29 incl. falle. Es wird folglich $D=23+19a-30p$,
woraus man leicht schliesst, dass $D$ der Rest sei, der entsteht, wenn man
$23+19a$ mit 30 dividirt, folglich $D=d$, oder die Oster-Gränze fällt auf den
$21+d^{\mathrm{ten}}$ März.

\SourceLeaf{81}%
III. Ostern selbst fällt nur auf den \textit{ersten} Sonntag \textit{nach} der Oster-Gränze,
also wenigstens einen, höchstens sieben Tage später als diese, mithin gewiss
nicht vor den $22+d^{\mathrm{ten}}$ März. Nimmt man also an, Ostern falle auf den
$22+d+E^{\mathrm{ten}}$ März, so liegt $E$ zwischen den Gränzen 0 und 6 incl., und muss
durch die Bedingung bestimmt werden, dass dieser Tag ein Sonntag sei. Diese
Bedingung lässt sich rein arithmetisch auf folgende Art ausdrücken: die Zwischenzeit
zwischen dem $22+d+E^{\mathrm{ten}}$ März des vorgegebenen Jahres und irgend
einem bestimmten Sonntage muss eine durch 7 theilbare Zahl von Tagen (eine
volle Anzahl Wochen) ausmachen. Man muss also einen bestimmten Sonntag
annehmen; ich wähle dazu den 21. März 1700. Nennt man nun die Zahl des
vorgegebenen Jahres $A$, und $i$ die Anzahl der zwischen 1700 und dem Jahre $A$
enthaltenen Schaltjahre, dieses, wenn es eines ist, eingeschlossen, so wird $i$ zugleich
die Anzahl der zwischen den 21. März 1700 und Ostern des Jahres $A$ eingefallenen
Schalttage sein, und die Anzahl \textit{aller} Tage vom 21. März 1700 bis
zum $22+d+E^{\mathrm{ten}}$ März des Jahres $A$
\[
=1+d+E+i+365(A-1700)
\]
Eben so leicht erhellt, dass zwischen 1700 und 1799 sein werde
\[
i=\tfrac14(A-b-1700)
\]
zwischen 1800 und 1899 hingegen
\[
i=\tfrac14(A-b-1700)-1
\]
Zur Bestimmung von $E$ hat man also die Bedingung, dass
\[
1+d+E+365(A-1700)+\tfrac14(A-b-1700)
\]
oder
\[
d+E+365(A-1700)+\tfrac14(A-b-1700)
\]
durch 7 theilbar sein müsse, je nachdem das Jahr zwischen 1700 und 1799 oder
zwischen 1800 und 1899 fällt. Es muss also auch eine durch 7 theilbare Zahl
herauskommen, wenn man ein Vielfaches von 7 zu jener addirt, oder davon abzieht,
oder auch jene von einem Vielfachen von 7 abzieht. Ich addire zuvörderst,
um den Bruch wegzuschaffen, $\frac74(A-b-1700)$, welches, wie man leicht
sieht, durch 7 theilbar ist; daraus erhalte ich
